\section{The ReGAE Model Family}
\label{sec:method}

A conditional Regenerative AutoEncoder (ReGAE) combines a compressed representation with a denoising operator that can be applied repeatedly to a mutable state. The representation specifies what to reconstruct, while the state supplies the current reconstruction, partial observation, or noise initialization. An unconditional member omits the encoder and its conditioning, retaining the recurrent denoiser. We describe these variants independently of the dimensionality, network widths, loss, and sampling choices of the Pandora's Cluster case study.

\subsection{State, encoding, and conditional update}

Let $\mathcal{X}$ be a space of sample tensors and $\mathcal{Z}$ a latent tensor space. An encoder $E_\phi:\mathcal{X}\to\mathcal{Z}$ compresses a sample $x$ into $z=E_\phi(x)$. Neither space has a prescribed dimensionality: the latent may be a vector or a structured tensor, with a bottleneck chosen for the task. Compression here describes the learned representation; its storage cost also depends on numerical precision and model parameters.

The conditional denoiser has the interface
\begin{equation}
  D_\theta:\mathcal{X}\times\mathcal{Z}\to\mathcal{X},
  \qquad x_{t+1}=D_\theta(x_t,z).
  \label{eq:family-recurrence}
\end{equation}
For a chosen configuration, its state input and output have the same dimensions, allowing the output to become the next input without a separate encoding step. The same construction accommodates one-, two-, or three-dimensional state layouts, using suitable operators for sequences, images, or volumetric tensors. This is an architectural interface, rather than a claim that all such domains have been evaluated here.

The latent remains fixed during a conditional trajectory. Parameters are also shared across updates: iteration changes the state, not the learned weights. Conditioning can enter through feature modulation or another compatible mechanism. The encoder is needed to obtain a code from an observed sample; repeated denoising subsequently uses that code directly.

\subsection{Learning a regenerative denoiser}

Training presents a corrupted version of a clean sample and asks the denoiser to recover the sample given its encoding. Let $C_\alpha(x,\epsilon)$ denote a corruption operator, with severity $\alpha$ and random variable $\epsilon$. One useful choice is the mixture
\begin{equation}
  C_\alpha(x,\epsilon)=(1-\alpha)x+\alpha\epsilon,
  \qquad 0\leq\alpha\leq1.
  \label{eq:family-corruption}
\end{equation}
The noise distribution, severity distribution, and any spatial masking are configuration choices. Including both weak and strong corruption connects refinement of an existing state with reconstruction from noise. The conditional training objective is
\begin{equation}
  \mathcal{L}_{\mathrm{cond}}=
  \mathbb{E}_{x,\alpha,\epsilon}\left[
    \ell\bigl(D_\theta(C_\alpha(x,\epsilon),E_\phi(x)),x\bigr)
  \right].
  \label{eq:family-conditional-loss}
\end{equation}
The reconstruction criterion $\ell$ may be an L1 loss, squared L2 loss, or feature-based \emph{perceptive} loss, chosen according to the data and desired fidelity. A feature loss compares representations of the prediction and target; it does not define the ReGAE family. Likewise, a particular noise schedule or optimizer is an implementation choice.

Equation~\eqref{eq:family-conditional-loss} trains an individual restoration step. Repeated application at inference is a separate use of that learned operator. Good one-step reconstruction does not alone establish convergence of every trajectory, preservation of all details, or recovery from arbitrary perturbations. Regeneration must therefore be assessed for the corruption types and update horizons of interest.

\subsection{Generation and regeneration}

Conditional generation starts from a noise state $x_0\sim p_0$ and applies Eq.~\eqref{eq:family-recurrence} for a chosen number of updates. When $z=E_\phi(x)$, this constructs a reconstruction associated with the encoded sample. Generating new samples through new codes additionally requires a way to choose those codes; the conditional autoencoder alone does not specify their prior distribution.

Regeneration uses the same update after an external perturbation. If $\widehat{x}$ is a current reconstruction and $P$ describes damage or corruption, set
\begin{equation}
  y_0=P(\widehat{x}),\qquad
  y_{t+1}=D_\theta(y_t,z).
  \label{eq:family-regeneration}
\end{equation}
The unchanged code supplies a reference while the denoiser acts on the altered state. Noise injection or partial-state refresh between updates can also be used as sampling choices. These interventions change the trajectory without requiring a new encoder evaluation or parameter update. A stable reconstruction may behave as an approximate fixed point, but fixed-point existence, attraction, and universal convergence are not assumed by the definition.

\subsection{Unconditional variant and architecture}

An unconditional ReGAE omits the encoder and external latent code, retaining the state-to-state denoiser:
\begin{equation}
  \begin{aligned}
    x_{t+1}&=D_\theta(x_t),\\
    \mathcal{L}_{\mathrm{uncond}}&=
    \mathbb{E}\left[\ell(D_\theta(C_\alpha(x,\epsilon)),x)\right].
  \end{aligned}
  \label{eq:family-unconditional}
\end{equation}
Here \emph{unconditional} means no separate conditioning code; the output still depends on its input state. Starting from noise gives a generative process shaped by the learned data regularities. Starting from an existing state instead uses its content as context. Without an external code, there is no independent encoded reference that identifies a particular original sample during recovery.

Our default denoiser uses a configurable U-Net-based architecture. Its depth, widths, normalization, attention, positional features, and conditioning mechanism can be selected for the problem while preserving matching state-input and state-output dimensions. Multiplicative feature gating is one optional component. For example, the gate called \emph{gravity} can apply $h\odot\gamma(g)$ to learned features, with a bounded or exponential $\gamma$. Such a gate changes feature gain; it is neither a required component of ReGAE nor a physical model of gravity. The precise choices used in the case study are documented separately.

\subsection{Optional hierarchy and predictive targets}

A hierarchy can apply the family to representations learned by another ReGAE. A full conditional first stage encodes samples or local units and provides a decoder for each code. Their codes can then be assembled into a latent field $Z$. An unconditional second-stage ReGAE, called the \emph{Gencoder}, learns a state update $Z_{t+1}=G_\psi(Z_t)$. Decoding its generated codes through the first stage maps latent organization back into the original data space. The two stages have separate learned parameters; a latent hierarchy is optional, rather than part of the definition of every ReGAE.

The restoration target can also be generalized to a same-shaped paired target. For paired states $(u,v)$ and optional conditioning $c$, train
\begin{equation}
  \mathcal{L}_{\mathrm{pair}}=
  \mathbb{E}\left[
    \ell\bigl(D_\theta(C_\alpha(u,\epsilon),c),v\bigr)
  \right],
  \label{eq:family-paired-target}
\end{equation}
omitting $c$ for an unconditional instance. Ordinary restoration takes $v=u$. A predictive extension instead uses a related displaced or subsequent state as $v$, so repeated updates may produce progression rather than convergence to one reconstruction. The Pandora experiment applies this extension to shifted latent fields. Its patch sizes, offset grids, overlap rules, and video sampler instantiate that particular hierarchy; they do not restrict the general ReGAE family.
